\chapter{Motivation und Zielsetzung}
\pagenumbering{arabic}
\textbf{Johanna Hoppe}


Entzündungen des Gehirns (Enzephalitis) werden meist durch Bakterien
oder Viren ausgelöst. Seltener ist eine Fehlreaktion des Immunsystems
der Grund.

Die Autoimmune Enzephalitis ist eine neurologische Erkrankung, die durch
Autoantikörper gegen neuronale Oberflächenproteine ausgelöst wird. Diese
fehlgeleitete Immunreaktion des Körpers führt zu einer beeinträchtigten
Signalweiterleitung am Axon und einer gestörten synaptischen
Übertragung. Dies äußert sich durch Symptome wie Gedächtnisstörungen, psychiatrischen Auffälligkeiten, oder Wesensverändernung. Betroffene leiden häufig
an verzögerter Muskelentspannung nach willkürlicher Muskelanspannung,
epileptischen Anfällen sowie Schlafstörungen und kognitiven
Defiziten.\footcite{VanSonderenNatRev2016}

Der Fokus der vorliegenden Arbeit liegt auf der autoimmunen
Caspr2-Enzephalitis, bei der Autoantikörper gegen das neuronale Protein
Contactin-associated protein 2 (Caspr2) gebildet werden. Caspr2 bildet mit Kaliumionenkanälen und
weiteren Proteinen im Nervensystem Komplexe, welche unter anderem für
die Lokalisation der Kaliumionenkanäle relevant sind. Bekannt ist
bereits, dass Caspr2 für die funktionierende Signalweiterleitung durch
das Axon von Bedeutung ist. Studien konnten Caspr2 2008 ebenfalls an Synapsen, den Verbindungsstellen zwischen zwei Nervenzellen, nachweisen. Der Fokus dieser Arbeit liegt daher auf der
Untersuchung von Caspr2 und der Wirkung von Autoantikörpern gegen
dieses Protein an den Synapsen des zentralen Nervensystems. In
Zusammenarbeit mit Dr.~Josefine Sell vom Universitätsklinikum Jena
werden zwei Teiluntersuchungen
durchgeführt.

Zum einen soll herausgefunden werden, ob Caspr2 eher an erregenden (exzitatorischen) oder
hemmenden (inhibitorischen) Synapsen lokalisiert ist. Dazu werden an
Hippocampuszellkulturen von Mäusen Doppelfärbungen kommerzieller
Caspr2-Antikörper mit je einem der synaptischen Marker VGlut1 (für
exzitatorische Synapsen) und VGat (für inhibitorische Synapsen)
durchgeführt. Diese werden mithilfe hochauflösender
Fluoreszenzmikroskopie analysiert und ausgewertet.

Im zweiten Teil wird untersucht, ob Antikörper gegen Caspr2 die
Interaktion zwischen Caspr2 und seinem Bindungspartner Contactin-2
stören. Dafür werden lebende neuronale Zellkulturen mit aus Patienten gewonnenen 
Caspr2-Antikörpern behandelt. Anschließend erfolgt eine Doppelfärbung
von Caspr2 und Contactin-2, um mögliche Veränderungen in der
Kolokalisation der beiden Proteine festzustellen.

Ziel dieser Seminarfacharbeit ist es, durch die Untersuchungen dazu
beizutragen, ein besseres Verständnis für die Caspr2-assoziierte
Enzephalitis zu entwickeln. Denn nur wenn die zugrundeliegenden Vorgänge
im Gehirn genau verstanden werden, kann die Diagnostik und Behandlung
der Krankheit weiter optimiert werden.